\documentclass[]{article}
\usepackage{lmodern}
\usepackage{amssymb,amsmath}
\usepackage{ifxetex,ifluatex}
\usepackage{fixltx2e} % provides \textsubscript
\ifnum 0\ifxetex 1\fi\ifluatex 1\fi=0 % if pdftex
  \usepackage[T1]{fontenc}
  \usepackage[utf8]{inputenc}
\else % if luatex or xelatex
  \ifxetex
    \usepackage{mathspec}
  \else
    \usepackage{fontspec}
  \fi
  \defaultfontfeatures{Ligatures=TeX,Scale=MatchLowercase}
\fi
% use upquote if available, for straight quotes in verbatim environments
\IfFileExists{upquote.sty}{\usepackage{upquote}}{}
% use microtype if available
\IfFileExists{microtype.sty}{%
\usepackage{microtype}
\UseMicrotypeSet[protrusion]{basicmath} % disable protrusion for tt fonts
}{}
\usepackage[margin=1in]{geometry}
\usepackage[unicode=true]{hyperref}
\hypersetup{
            pdfborder={0 0 0},
            breaklinks=true}
\urlstyle{same}  % don't use monospace font for urls
\usepackage{color}
\usepackage{fancyvrb}
\newcommand{\VerbBar}{|}
\newcommand{\VERB}{\Verb[commandchars=\\\{\}]}
\DefineVerbatimEnvironment{Highlighting}{Verbatim}{commandchars=\\\{\}}
% Add ',fontsize=\small' for more characters per line
\usepackage{framed}
\definecolor{shadecolor}{RGB}{248,248,248}
\newenvironment{Shaded}{\begin{snugshade}}{\end{snugshade}}
\newcommand{\KeywordTok}[1]{\textcolor[rgb]{0.13,0.29,0.53}{\textbf{{#1}}}}
\newcommand{\DataTypeTok}[1]{\textcolor[rgb]{0.13,0.29,0.53}{{#1}}}
\newcommand{\DecValTok}[1]{\textcolor[rgb]{0.00,0.00,0.81}{{#1}}}
\newcommand{\BaseNTok}[1]{\textcolor[rgb]{0.00,0.00,0.81}{{#1}}}
\newcommand{\FloatTok}[1]{\textcolor[rgb]{0.00,0.00,0.81}{{#1}}}
\newcommand{\ConstantTok}[1]{\textcolor[rgb]{0.00,0.00,0.00}{{#1}}}
\newcommand{\CharTok}[1]{\textcolor[rgb]{0.31,0.60,0.02}{{#1}}}
\newcommand{\SpecialCharTok}[1]{\textcolor[rgb]{0.00,0.00,0.00}{{#1}}}
\newcommand{\StringTok}[1]{\textcolor[rgb]{0.31,0.60,0.02}{{#1}}}
\newcommand{\VerbatimStringTok}[1]{\textcolor[rgb]{0.31,0.60,0.02}{{#1}}}
\newcommand{\SpecialStringTok}[1]{\textcolor[rgb]{0.31,0.60,0.02}{{#1}}}
\newcommand{\ImportTok}[1]{{#1}}
\newcommand{\CommentTok}[1]{\textcolor[rgb]{0.56,0.35,0.01}{\textit{{#1}}}}
\newcommand{\DocumentationTok}[1]{\textcolor[rgb]{0.56,0.35,0.01}{\textbf{\textit{{#1}}}}}
\newcommand{\AnnotationTok}[1]{\textcolor[rgb]{0.56,0.35,0.01}{\textbf{\textit{{#1}}}}}
\newcommand{\CommentVarTok}[1]{\textcolor[rgb]{0.56,0.35,0.01}{\textbf{\textit{{#1}}}}}
\newcommand{\OtherTok}[1]{\textcolor[rgb]{0.56,0.35,0.01}{{#1}}}
\newcommand{\FunctionTok}[1]{\textcolor[rgb]{0.00,0.00,0.00}{{#1}}}
\newcommand{\VariableTok}[1]{\textcolor[rgb]{0.00,0.00,0.00}{{#1}}}
\newcommand{\ControlFlowTok}[1]{\textcolor[rgb]{0.13,0.29,0.53}{\textbf{{#1}}}}
\newcommand{\OperatorTok}[1]{\textcolor[rgb]{0.81,0.36,0.00}{\textbf{{#1}}}}
\newcommand{\BuiltInTok}[1]{{#1}}
\newcommand{\ExtensionTok}[1]{{#1}}
\newcommand{\PreprocessorTok}[1]{\textcolor[rgb]{0.56,0.35,0.01}{\textit{{#1}}}}
\newcommand{\AttributeTok}[1]{\textcolor[rgb]{0.77,0.63,0.00}{{#1}}}
\newcommand{\RegionMarkerTok}[1]{{#1}}
\newcommand{\InformationTok}[1]{\textcolor[rgb]{0.56,0.35,0.01}{\textbf{\textit{{#1}}}}}
\newcommand{\WarningTok}[1]{\textcolor[rgb]{0.56,0.35,0.01}{\textbf{\textit{{#1}}}}}
\newcommand{\AlertTok}[1]{\textcolor[rgb]{0.94,0.16,0.16}{{#1}}}
\newcommand{\ErrorTok}[1]{\textcolor[rgb]{0.64,0.00,0.00}{\textbf{{#1}}}}
\newcommand{\NormalTok}[1]{{#1}}
\usepackage{longtable,booktabs}
\usepackage{graphicx,grffile}
\makeatletter
\def\maxwidth{\ifdim\Gin@nat@width>\linewidth\linewidth\else\Gin@nat@width\fi}
\def\maxheight{\ifdim\Gin@nat@height>\textheight\textheight\else\Gin@nat@height\fi}
\makeatother
% Scale images if necessary, so that they will not overflow the page
% margins by default, and it is still possible to overwrite the defaults
% using explicit options in \includegraphics[width, height, ...]{}
\setkeys{Gin}{width=\maxwidth,height=\maxheight,keepaspectratio}
\IfFileExists{parskip.sty}{%
\usepackage{parskip}
}{% else
\setlength{\parindent}{0pt}
\setlength{\parskip}{6pt plus 2pt minus 1pt}
}
\setlength{\emergencystretch}{3em}  % prevent overfull lines
\providecommand{\tightlist}{%
  \setlength{\itemsep}{0pt}\setlength{\parskip}{0pt}}
\setcounter{secnumdepth}{0}
% Redefines (sub)paragraphs to behave more like sections
\ifx\paragraph\undefined\else
\let\oldparagraph\paragraph
\renewcommand{\paragraph}[1]{\oldparagraph{#1}\mbox{}}
\fi
\ifx\subparagraph\undefined\else
\let\oldsubparagraph\subparagraph
\renewcommand{\subparagraph}[1]{\oldsubparagraph{#1}\mbox{}}
\fi

\author{}
\date{\vspace{-2.5em}}

\begin{document}

\section{X2Chess Architecture Manual}\label{x2chess-architecture-manual}

Last updated: 2026-04-04

\subsection{1) Goals}\label{goals}

\begin{itemize}
\tightlist
\item
  Make the browser fat client the primary/full product runtime.
\item
  Allow remote launch (server-hosted assets) while preserving local
  client execution.
\item
  Combine local resource usage with remote resource integration in one
  client architecture.
\item
  Provide an app-profile variant as a reduced version of the full fat
  client.
\item
  Keep architecture decisions explicit and traceable.
\end{itemize}

Architecture decisions are maintained in
\texttt{doc/architecture-adr-manual.qmd}.

\subsection{2) System Overview}\label{system-overview}

\subsection{\texorpdfstring{2.1 Frontend
(\texttt{frontend/})}{2.1 Frontend (frontend/)}}\label{frontend-frontend}

\begin{itemize}
\tightlist
\item
  Technology: Vite + React 18 + TypeScript (ES modules).
\item
  Main entrypoint: \texttt{frontend/src/main.tsx}.
\item
  React state: single \texttt{useReducer} store (\texttt{AppStoreState})
  owned by \texttt{AppProvider}; all mutations via
  \texttt{dispatch(action)}.
\item
  Service layer: \texttt{useAppStartup} hook initialises all pure-logic
  services on mount, wraps them in a \texttt{SessionOrchestrator}, and
  exposes stable callbacks via \texttt{ServiceContext}; components
  consume them with \texttt{useServiceContext()}. Each service owns its
  internal state as closure variables; services communicate state
  changes to React via typed callbacks that dispatch fine-grained
  actions (see §7.3).
\item
  Responsibilities:
\item
  Render the Game-Viewer (board + PGN editor UI) as React components.
\item
  Parse PGN into an internal model and serialize back to text.
\item
  Keep replay/navigation state in React \texttt{useReducer}.
\item
  Manage chess move playback and UI interactions via
  \texttt{ServiceContext} callbacks.
\item
  Manage multiple open game sessions in parallel tabs.
\item
  Route load/save through a uniform source adapter contract.
\end{itemize}

\subsection{\texorpdfstring{2.2 Backend
(\texttt{backend/})}{2.2 Backend (backend/)}}\label{backend-backend}

\begin{itemize}
\tightlist
\item
  Technology: FastAPI.
\item
  Scope in current architecture:
\item
  Optional static serving and optional future remote services.
\item
  Not part of client-local resource loading/saving.
\item
  Constraint:
\item
  Local runtime resources are owned and handled only by the client
  runtime.
\end{itemize}

\subsection{\texorpdfstring{2.3 App Profile Shell
(\texttt{frontend/src-tauri/})}{2.3 App Profile Shell (frontend/src-tauri/)}}\label{app-profile-shell-frontendsrc-tauri}

\begin{itemize}
\tightlist
\item
  Technology: Tauri 2 (Rust host + web frontend).
\item
  Configuration: \texttt{frontend/src-tauri/tauri.conf.json}.
\item
  Runtime entrypoint: \texttt{frontend/src-tauri/src/main.rs}.
\item
  Capability model:
\item
  \texttt{frontend/src-tauri/capabilities/default.json} with
  \texttt{core:default}, \texttt{core:menu:default}.
\item
  Build wiring:
\item
  Dev URL \texttt{http://localhost:5287}.
\item
  Production frontend assets from \texttt{frontend/dist}.
\end{itemize}

\subsubsection{Desktop menu bar}\label{desktop-menu-bar}

The native desktop menu bar is defined declaratively in a single
TypeScript file:

\begin{itemize}
\item
  \textbf{Definition:}
  \texttt{frontend/src/app\_shell/menu\_definition.ts} --- a plain,
  framework-free data structure (\texttt{APP\_MENU\_DEFINITION}) that
  lists every top-level menu and its items (actions, predefined OS
  items, separators, submenus). Adding, removing, or reordering menu
  items only requires editing this file.
\item
  \textbf{Action ids} (\texttt{MenuActionId}) name the application-level
  operations (e.g. \texttt{"file.open-resource"},
  \texttt{"file.new-database"}). Each id is handled by a matching entry
  in the \texttt{MENU\_ACTIONS} map inside \texttt{useTauriMenu.ts}.
\item
  \textbf{Build hook:} \texttt{frontend/src/hooks/useTauriMenu.ts} --- a
  React hook called once by \texttt{useAppStartup}. It reads
  \texttt{APP\_MENU\_DEFINITION}, constructs the Tauri 2 \texttt{Menu}
  via \texttt{@tauri-apps/api/menu}, wires action callbacks to
  \texttt{AppStartupServices}, and installs the result with
  \texttt{menu.setAsAppMenu()}. The hook is a no-op in browser runtime
  (\texttt{isTauriRuntime()} guard).
\end{itemize}

\textbf{To add a new menu action:} 1. Add a new id to
\texttt{MenuActionId} in \texttt{menu\_definition.ts}. 2. Insert a
\texttt{\{\ kind:\ "action",\ id:\ "\ldots{}",\ label:\ "\ldots{}",\ accelerator?:\ "\ldots{}"\ \}}
node at the desired position in \texttt{APP\_MENU\_DEFINITION}. 3. Add
the handler
\texttt{"new.id":\ (services)\ =\textgreater{}\ \{\ \ldots{}\ \}} in
\texttt{MENU\_ACTIONS} inside \texttt{useTauriMenu.ts}.

\subsection{\texorpdfstring{2.4 Resource Library
(\texttt{resource/})}{2.4 Resource Library (resource/)}}\label{resource-library-resource}

\begin{itemize}
\tightlist
\item
  Technology: TypeScript library consumed by frontend runtime
  boundaries.
\item
  Scope in current architecture:
\item
  Defines canonical PGN resource contracts (\texttt{PgnResourceRef},
  \texttt{PgnGameRef}, \texttt{PgnGameEntry}).
\item
  Provides canonical kind model: \texttt{file}, \texttt{directory},
  \texttt{db}.
\item
  Owns metadata schema/constants and extraction contracts.
\item
  Owns database schema and migrations under \texttt{resource/database/}.
\item
  Boundary constraints:
\item
  Frontend modules consume \texttt{resource/client} and
  \texttt{resource/domain} contracts.
\item
  Frontend must not own DB schema/migration internals.
\end{itemize}

\subsection{3) Build and Delivery Paths}\label{build-and-delivery-paths}

Environment and toolchain setup is documented in
\texttt{doc/setup-manual.qmd}.

\subsection{3.1 Full browser fat-client
profile}\label{full-browser-fat-client-profile}

\begin{itemize}
\tightlist
\item
  Dev (local): \texttt{cd\ frontend\ \&\&\ npm\ run\ dev}
\item
  Build (deployable from remote host or local static host):
  \texttt{cd\ frontend\ \&\&\ npm\ run\ build}
\end{itemize}

\subsection{3.2 Reduced app profile
(Tauri)}\label{reduced-app-profile-tauri}

\begin{itemize}
\tightlist
\item
  Dev shell + hot reload:
  \texttt{cd\ frontend\ \&\&\ npm\ run\ desktop:dev}
\item
  Packaged app profile:
  \texttt{cd\ frontend\ \&\&\ npm\ run\ desktop:pack}
\item
  Artifact directory: \texttt{frontend/src-tauri/target/release/bundle/}
\end{itemize}

\subsection{3.3 Runtime modes and
Developer-Tools}\label{runtime-modes-and-developer-tools}

\begin{itemize}
\tightlist
\item
  Build mode is injected as \texttt{\_\_X2CHESS\_MODE\_\_} (\texttt{DEV}
  or \texttt{PROD}).
\item
  Default behavior:
\item
  \texttt{DEV}: Developer-Tools starts ON.
\item
  \texttt{PROD}: Developer-Tools starts OFF.
\item
  Developer-Tools can be toggled in the app menu and persists per user.
\item
  Developer-Tools exposes a bottom fixed dock with tabs:
\item
  \texttt{AST}
\item
  \texttt{DOM}
\item
  \texttt{Raw\ PGN}
\item
  The dock is user-opened and user-closed; enabling Developer-Tools does
  not force the dock open.
\item
  Dock views always reflect the same in-memory game/model state as the
  main Game-Viewer.
\item
  DEV-only convenience:
\item
  Automatic preload attempt from \texttt{run/DEV/games} (when
  available).
\end{itemize}

\subsection{3.4 Game sessions and source
adapters}\label{game-sessions-and-source-adapters}

\begin{itemize}
\tightlist
\item
  Open games are modeled as in-memory \texttt{GameSession} entries:
\item
  \texttt{sessionId}
\item
  \texttt{sourceRef}
\item
  \texttt{revisionToken}
\item
  snapshot of editor/board runtime
\item
  dirty/save mode metadata
\item
  Session identity is independent from source identity; source
  references are attached metadata.
\item
  Source interactions are routed through canonical resource boundaries:
\item
  \texttt{file} kind: single file containing multiple games
\item
  \texttt{directory} kind: folder of one-game-per-file entries
\item
  \texttt{db} kind: database-backed resources (schema/migrations owned
  in \texttt{resource/database/}; implementation may be deferred by
  profile)
\item
  Supported open flows are:
\item
  Drag/drop PGN text/files
\item
  Paste PGN text
\item
  Source-root preload (DEV convenience)
\item
  Save path is uniform:
\item
  Active session + source reference + revision token -\textgreater{}
  source adapter \texttt{save(...)}.
\item
  Save policy:
\item
  Default mode is autosave.
\item
  User can switch active-session mode between autosave and manual from
  the app menu.
\item
  Naming convention:
\item
  The conjunction of board and text editor is called
  \textbf{Game-Viewer}.
\end{itemize}

\subsection{3.5 App-Version and Data-Migration
Strategy}\label{app-version-and-data-migration-strategy}

All persistent data in X2Chess is versioned so that older instances can
be detected and upgraded automatically on startup. Three layers are
covered:

\subsubsection{3.5.1 App version constant}\label{app-version-constant}

\begin{itemize}
\tightlist
\item
  \textbf{Source of truth:} \texttt{frontend/src-tauri/tauri.conf.json}
  → injected as \texttt{\_\_X2CHESS\_APP\_VERSION\_\_} by Vite.
\item
  \textbf{Access point:} \texttt{frontend/src/runtime/app\_version.ts}
  exports \texttt{CURRENT\_APP\_VERSION:\ string}, \texttt{parseSemver},
  and \texttt{isNewerVersion}. All code that needs the version reads
  from this module; no other file duplicates the inline
  \texttt{typeof\ \_\_X2CHESS\_APP\_VERSION\_\_} guard.
\end{itemize}

\subsubsection{3.5.2 SQLite database schema
migrations}\label{sqlite-database-schema-migrations}

\begin{itemize}
\tightlist
\item
  \textbf{Owner:} \texttt{resource/database/migrations/} (runner + steps
  + manifest).
\item
  \textbf{Mechanism:} \texttt{schema\_version} table records applied
  versions. \texttt{runMigrations(db)} is idempotent --- safe to call on
  every DB open; it skips already-applied steps.
\item
  \textbf{Naming convention:} migration files are prefixed
  \texttt{NNNN\_} (e.g. \texttt{0001\_init.ts}).
\item
  \textbf{Rule:} Never modify a migration step after it has been
  shipped. Add a new step instead.
\end{itemize}

\subsubsection{3.5.3 localStorage versioned
stores}\label{localstorage-versioned-stores}

\begin{itemize}
\tightlist
\item
  \textbf{Primitive:} \texttt{frontend/src/storage/versioned\_store.ts}
  --- \texttt{createVersionedStore\textless{}T\textgreater{}(config)}.
\item
  \textbf{Envelope format:} every managed key stores
  \texttt{\{\ "v":\ \textless{}number\textgreater{},\ "data":\ \textless{}T\textgreater{}\ \}}
  JSON.
\item
  \textbf{On read:}
\item
  Absent key → return default (no write-back).
\item
  Current version → return data directly.
\item
  Older version → apply migration chain in-place; write upgraded
  envelope back.
\item
  Future version (data from a newer build) → return default
  \emph{without overwriting}.
\item
  Corrupt JSON → reset key to default; log a warning.
\item
  \textbf{Adding a store:} call \texttt{createVersionedStore} with
  \texttt{version:\ 1} and \texttt{migrations:\ {[}{]}}. Increment
  \texttt{version} and add a migration step whenever the stored shape
  changes incompatibly.
\item
  \textbf{Migration step indexing:} \texttt{migrations{[}i{]}} upgrades
  data from version \texttt{i} to version \texttt{i+1}. Array length
  must equal \texttt{version\ -\ 1}. A step returning \texttt{null}
  signals unrecoverable data and triggers a reset to
  \texttt{defaultValue}.
\end{itemize}

\paragraph{Managed stores (version 1
baseline)}\label{managed-stores-version-1-baseline}

\begin{longtable}[]{@{}lll@{}}
\toprule
\begin{minipage}[b]{0.05\columnwidth}\raggedright\strut
Store\strut
\end{minipage} & \begin{minipage}[b]{0.05\columnwidth}\raggedright\strut
Key\strut
\end{minipage} & \begin{minipage}[b]{0.05\columnwidth}\raggedright\strut
File\strut
\end{minipage}\tabularnewline
\midrule
\endhead
\begin{minipage}[t]{0.05\columnwidth}\raggedright\strut
Shell prefs (compound)\strut
\end{minipage} & \begin{minipage}[t]{0.05\columnwidth}\raggedright\strut
\texttt{x2chess.shellPrefs}\strut
\end{minipage} & \begin{minipage}[t]{0.05\columnwidth}\raggedright\strut
\texttt{runtime/shell\_prefs\_store.ts}\strut
\end{minipage}\tabularnewline
\begin{minipage}[t]{0.05\columnwidth}\raggedright\strut
Shape / board decoration prefs\strut
\end{minipage} & \begin{minipage}[t]{0.05\columnwidth}\raggedright\strut
\texttt{x2chess.shapePrefs.v1}\strut
\end{minipage} & \begin{minipage}[t]{0.05\columnwidth}\raggedright\strut
\texttt{runtime/shape\_prefs.ts}\strut
\end{minipage}\tabularnewline
\begin{minipage}[t]{0.05\columnwidth}\raggedright\strut
Resource viewer column prefs\strut
\end{minipage} & \begin{minipage}[t]{0.05\columnwidth}\raggedright\strut
\texttt{x2chess.resourceViewerColumnPrefs.v1}\strut
\end{minipage} & \begin{minipage}[t]{0.05\columnwidth}\raggedright\strut
\texttt{resources\_viewer/viewer\_utils.ts}\strut
\end{minipage}\tabularnewline
\begin{minipage}[t]{0.05\columnwidth}\raggedright\strut
Metadata schemas\strut
\end{minipage} & \begin{minipage}[t]{0.05\columnwidth}\raggedright\strut
\texttt{x2chess.metadata-schemas}\strut
\end{minipage} & \begin{minipage}[t]{0.05\columnwidth}\raggedright\strut
\texttt{resources\_viewer/schema\_storage.ts}\strut
\end{minipage}\tabularnewline
\begin{minipage}[t]{0.05\columnwidth}\raggedright\strut
External DB settings\strut
\end{minipage} & \begin{minipage}[t]{0.05\columnwidth}\raggedright\strut
\texttt{x2chess.ext-db-settings}\strut
\end{minipage} & \begin{minipage}[t]{0.05\columnwidth}\raggedright\strut
\texttt{runtime/ext\_database\_settings\_store.ts}\strut
\end{minipage}\tabularnewline
\begin{minipage}[t]{0.05\columnwidth}\raggedright\strut
Remote rules cache\strut
\end{minipage} & \begin{minipage}[t]{0.05\columnwidth}\raggedright\strut
\texttt{x2chess.webImportRules}\strut
\end{minipage} & \begin{minipage}[t]{0.05\columnwidth}\raggedright\strut
\texttt{runtime/remote\_rules\_store.ts}\strut
\end{minipage}\tabularnewline
\begin{minipage}[t]{0.05\columnwidth}\raggedright\strut
User web-import rules\strut
\end{minipage} & \begin{minipage}[t]{0.05\columnwidth}\raggedright\strut
\texttt{x2chess.webImport.userRules}\strut
\end{minipage} & \begin{minipage}[t]{0.05\columnwidth}\raggedright\strut
\texttt{resources/web\_import/user\_rules\_storage.ts}\strut
\end{minipage}\tabularnewline
\begin{minipage}[t]{0.05\columnwidth}\raggedright\strut
Training badges\strut
\end{minipage} & \begin{minipage}[t]{0.05\columnwidth}\raggedright\strut
\texttt{x2chess.training-badges}\strut
\end{minipage} & \begin{minipage}[t]{0.05\columnwidth}\raggedright\strut
\texttt{training/transcript\_storage.ts}\strut
\end{minipage}\tabularnewline
\begin{minipage}[t]{0.05\columnwidth}\raggedright\strut
Training sessions\strut
\end{minipage} & \begin{minipage}[t]{0.05\columnwidth}\raggedright\strut
\texttt{x2chess.training-sessions}\strut
\end{minipage} & \begin{minipage}[t]{0.05\columnwidth}\raggedright\strut
\texttt{training/transcript\_storage.ts}\strut
\end{minipage}\tabularnewline
\begin{minipage}[t]{0.05\columnwidth}\raggedright\strut
Curriculum plan\strut
\end{minipage} & \begin{minipage}[t]{0.05\columnwidth}\raggedright\strut
\texttt{x2chess.curriculum-plan}\strut
\end{minipage} & \begin{minipage}[t]{0.05\columnwidth}\raggedright\strut
\texttt{training/curriculum/curriculum\_storage.ts}\strut
\end{minipage}\tabularnewline
\begin{minipage}[t]{0.05\columnwidth}\raggedright\strut
Workspace snapshot\strut
\end{minipage} & \begin{minipage}[t]{0.05\columnwidth}\raggedright\strut
\texttt{x2chess.workspaceSnapshot}\strut
\end{minipage} & \begin{minipage}[t]{0.05\columnwidth}\raggedright\strut
\texttt{runtime/workspace\_snapshot\_store.ts}\strut
\end{minipage}\tabularnewline
\bottomrule
\end{longtable}

Dynamic per-tab group-by state
(\texttt{x2chess.groupby.\textless{}tabId\textgreater{}}) and the
visual-asset cache
(\texttt{chess-app:visual-asset:v2:\textless{}key\textgreater{}}) are
intentionally excluded --- they are ephemeral UI state and asset caches
respectively, not user configuration.

\paragraph{Physical storage location}\label{physical-storage-location}

X2Chess is a Tauri v2 desktop app (bundle ID \texttt{com.x2chess.app}).
Each platform's webview engine stores localStorage as opaque binary
files --- they are \textbf{not} standard SQLite databases and are not
intended to be edited directly.

\textbf{Default paths} (used when \texttt{X2CHESS\_WEBVIEW\_DATA} is not
set):

\begin{longtable}[]{@{}lll@{}}
\toprule
\begin{minipage}[b]{0.05\columnwidth}\raggedright\strut
Platform\strut
\end{minipage} & \begin{minipage}[b]{0.05\columnwidth}\raggedright\strut
Webview\strut
\end{minipage} & \begin{minipage}[b]{0.05\columnwidth}\raggedright\strut
Path\strut
\end{minipage}\tabularnewline
\midrule
\endhead
\begin{minipage}[t]{0.05\columnwidth}\raggedright\strut
macOS\strut
\end{minipage} & \begin{minipage}[t]{0.05\columnwidth}\raggedright\strut
WKWebView\strut
\end{minipage} & \begin{minipage}[t]{0.05\columnwidth}\raggedright\strut
\texttt{\textasciitilde{}/Library/WebKit/com.x2chess.app/WebsiteData/LocalStorage/}\strut
\end{minipage}\tabularnewline
\begin{minipage}[t]{0.05\columnwidth}\raggedright\strut
Windows\strut
\end{minipage} & \begin{minipage}[t]{0.05\columnwidth}\raggedright\strut
WebView2 (Chromium)\strut
\end{minipage} & \begin{minipage}[t]{0.05\columnwidth}\raggedright\strut
\texttt{\%APPDATA\%\textbackslash{}X2Chess\textbackslash{}EBWebView\textbackslash{}Default\textbackslash{}Local\ Storage\textbackslash{}leveldb\textbackslash{}}\strut
\end{minipage}\tabularnewline
\begin{minipage}[t]{0.05\columnwidth}\raggedright\strut
Linux\strut
\end{minipage} & \begin{minipage}[t]{0.05\columnwidth}\raggedright\strut
WebView2 (Chromium)\strut
\end{minipage} & \begin{minipage}[t]{0.05\columnwidth}\raggedright\strut
\texttt{\$HOME/.local/share/X2Chess/EBWebView/Default/Local\ Storage/leveldb\textbackslash{}}\strut
\end{minipage}\tabularnewline
\bottomrule
\end{longtable}

On \textbf{macOS} each origin maps to one file named
\texttt{\textless{}scheme\textgreater{}\_\textless{}host\textgreater{}\_\textless{}port\textgreater{}\_0.localstorage}
(WebKit binary format). The Tauri custom protocol origin is
\texttt{tauri://localhost}, so all \texttt{x2chess.*} keys share the
single file \texttt{tauri\_localhost\_0.localstorage}.

On \textbf{Windows / Linux} (WebView2) the storage is a LevelDB
directory; individual keys are not addressable as separate files.

\paragraph{\texorpdfstring{Redirecting storage via
\texttt{X2CHESS\_WEBVIEW\_DATA}}{Redirecting storage via X2CHESS\_WEBVIEW\_DATA}}\label{redirecting-storage-via-x2chess_webview_data}

The Tauri main window is created programmatically in \texttt{main.rs}
using \texttt{WebviewWindowBuilder::data\_directory()}. Setting the
environment variable \texttt{X2CHESS\_WEBVIEW\_DATA} to any directory
path redirects \textbf{all} webview storage (localStorage, IndexedDB,
cache, cookies) to that directory. The directory is created
automatically on first launch if it does not exist.

\begin{Shaded}
\begin{Highlighting}[]
\CommentTok{# Isolated dev profile — data written to run/DEV/webview-data/ (relative to repo root)}
\OtherTok{X2CHESS_WEBVIEW_DATA=}\NormalTok{../run/DEV/webview-data  }\KeywordTok{tauri} \NormalTok{dev}

\CommentTok{# Packaged as npm scripts:}
\KeywordTok{npm} \NormalTok{run desktop:dev:isolated    }\CommentTok{# isolated (recommended for development)}
\KeywordTok{npm} \NormalTok{run desktop:dev             }\CommentTok{# OS default path (shared with any installed PROD build)}
\end{Highlighting}
\end{Shaded}

The directory layout inside a redirected path mirrors the platform
default:

\begin{longtable}[]{@{}ll@{}}
\toprule
Platform & Sub-path under \texttt{X2CHESS\_WEBVIEW\_DATA}\tabularnewline
\midrule
\endhead
macOS &
\texttt{WebsiteData/LocalStorage/tauri\_localhost\_0.localstorage}\tabularnewline
Windows / Linux &
\texttt{EBWebView/Default/Local\ Storage/leveldb/}\tabularnewline
\bottomrule
\end{longtable}

\textbf{Full reset of isolated dev storage:}

\begin{Shaded}
\begin{Highlighting}[]
\KeywordTok{rm} \NormalTok{-rf run/DEV/webview-data   }\CommentTok{# directory is recreated on next launch}
\end{Highlighting}
\end{Shaded}

\begin{quote}
Note: The \texttt{x2chess.devPrefsMode} key documented in §4 (dev prefs
mode switch) lives in the same storage location. A full reset also
clears this flag, so the next launch re-initialises it to
\texttt{"defaults"}.
\end{quote}

\subsubsection{3.5.4 Legacy key consolidation (one-shot
migration)}\label{legacy-key-consolidation-one-shot-migration}

Individual shell-preference keys written by earlier builds
(\texttt{x2chess.sound}, \texttt{x2chess.moveDelayMs},
\texttt{x2chess.locale}, \texttt{x2chess.pgnLayout},
\texttt{x2chess.developerTools},
\texttt{x2chess.positionPreviewOnHover},
\texttt{x2chess.resourceViewerHeightPx},
\texttt{x2chess.boardColumnWidthPx}) are merged into the compound
\texttt{x2chess.shellPrefs} key on the \textbf{first startup after
upgrade} by \texttt{migrateLocalStorage()} in
\texttt{frontend/src/storage/migrate\_local\_storage.ts}. The function
is idempotent: if the compound key already exists it returns
immediately. Legacy keys are deleted after consolidation.

Similarly, \texttt{migrateRemoteRulesCache()} in
\texttt{runtime/remote\_rules\_store.ts} merges the legacy two-key
remote rules cache (\texttt{x2chess.webImportRules.v1} +
\texttt{x2chess.webImportRules.version}) into the compound
\texttt{x2chess.webImportRules} key.

\subsubsection{3.5.5 Startup sequence}\label{startup-sequence}

\begin{verbatim}
useAppStartup mount effect
  │
  ├─ [SYNC] migrateLocalStorage()         ← consolidate legacy shell-pref keys
  ├─ [SYNC] migrateRemoteRulesCache()     ← consolidate legacy rules-cache keys
  ├─ [SYNC] versioned store reads         ← auto-migrate on first access
  │
  └─ [ASYNC, inside Tauri gateway open]
       runMigrations(db)                  ← idempotent DB schema migration
\end{verbatim}

\subsubsection{3.5.6 Workspace snapshot}\label{workspace-snapshot}

The full workspace --- all open game sessions and all open resource
viewer tabs --- is persisted automatically so it can be restored on the
next launch.

\textbf{Store:} \texttt{x2chess.workspaceSnapshot} (versioned, managed
by \texttt{runtime/workspace\_snapshot\_store.ts}).

\textbf{Shape:} \texttt{WorkspaceSnapshot} contains:

\begin{itemize}
\tightlist
\item
  \texttt{sessions{[}{]}} --- one \texttt{SessionSnap} per open tab:
  session ID, title, full PGN text (including any unsaved edits),
  \texttt{sourceRef}, save mode, dirty state, current ply, selected move
  ID, and PGN layout mode.
\item
  \texttt{activeSessionId} --- the session that was focused when the
  snapshot was written.
\item
  \texttt{resourceTabs{[}{]}} --- one \texttt{ResourceTabSnap} per open
  resource viewer tab: kind, locator, title.
\item
  \texttt{activeResourceTabId} --- the active resource tab.
\end{itemize}

\textbf{Write path:} A \texttt{useEffect} in \texttt{useAppStartup}
watches \texttt{state.sessions}, \texttt{state.activeSessionId},
\texttt{state.resourceViewerTabSnapshots}, and
\texttt{state.activeResourceTabId}. Any change schedules a debounced
write (500 ms). When the timer fires, the current snapshots are
assembled by calling
\texttt{bundle.sessionStore.buildSessionSnapshots()} and
\texttt{bundle.resourceViewer.buildTabSnapshots()} and written
atomically to the store. Rapid successive changes collapse into a single
write.

\textbf{Restore path:} \texttt{useAppStartup} reads the snapshot
synchronously before opening any session. When
\texttt{snapshot.sessions} is non-empty it reconstructs each session
from the stored PGN text and restores the active session pointer and all
resource viewer tabs. When the snapshot is empty (first launch, or after
a reset) a blank default session is opened instead.

\textbf{Close-guard:} \texttt{AppShell} registers a
\texttt{beforeunload} handler that calls \texttt{e.preventDefault()} /
sets \texttt{e.returnValue} whenever any session has
\texttt{dirtyState\ ===\ "dirty"}. This causes the browser (or Tauri
webview) to show the platform's ``unsaved changes --- leave anyway?''
dialog before the window closes.

\subsubsection{3.5.7 Version-gating
(future)}\label{version-gating-future}

Neither localStorage nor DB migrations currently require app-version
gating (all steps are additive). If a future migration must only run on
a specific version transition, use \texttt{CURRENT\_APP\_VERSION} /
\texttt{parseSemver} from \texttt{runtime/app\_version.ts} as the guard
inside the migration step or \texttt{migrateLocalStorage}.

\subsection{4) Configuration}\label{configuration}

Configuration in X2Chess splits into two distinct categories:
\textbf{build-time configuration} embedded at compile time and
\textbf{user preferences} persisted at runtime.

\subsubsection{4.1 Build-time
configuration}\label{build-time-configuration}

Build-time values are fixed when the application is compiled. They
cannot be changed at runtime.

\textbf{Source of truth:} \texttt{frontend/src-tauri/tauri.conf.json}

This file owns:

\begin{itemize}
\tightlist
\item
  \texttt{version} --- the application semver string (canonical source
  for \texttt{CURRENT\_APP\_VERSION}; see §3.5.1).
\item
  \texttt{productName}, \texttt{identifier} --- product and bundle
  identity.
\item
  Window dimensions and minimum constraints (\texttt{width},
  \texttt{height}, \texttt{minWidth}, \texttt{minHeight}).
\item
  Updater endpoint and signing public key.
\end{itemize}

\textbf{Vite build globals}

Vite reads \texttt{tauri.conf.json} at build time
(\texttt{frontend/vite.config.ts}) and injects three compile-time
constants into every bundle:

\begin{longtable}[]{@{}lll@{}}
\toprule
\begin{minipage}[b]{0.05\columnwidth}\raggedright\strut
Global\strut
\end{minipage} & \begin{minipage}[b]{0.05\columnwidth}\raggedright\strut
Value\strut
\end{minipage} & \begin{minipage}[b]{0.05\columnwidth}\raggedright\strut
Usage\strut
\end{minipage}\tabularnewline
\midrule
\endhead
\begin{minipage}[t]{0.05\columnwidth}\raggedright\strut
\texttt{\_\_X2CHESS\_MODE\_\_}\strut
\end{minipage} & \begin{minipage}[t]{0.05\columnwidth}\raggedright\strut
\texttt{"DEV"} or \texttt{"PROD"}\strut
\end{minipage} & \begin{minipage}[t]{0.05\columnwidth}\raggedright\strut
Controls developer-tools default; consumed by
\texttt{runtime/bootstrap\_prefs.ts}.\strut
\end{minipage}\tabularnewline
\begin{minipage}[t]{0.05\columnwidth}\raggedright\strut
\texttt{\_\_X2CHESS\_APP\_VERSION\_\_}\strut
\end{minipage} & \begin{minipage}[t]{0.05\columnwidth}\raggedright\strut
Semver string from \texttt{tauri.conf.json}\strut
\end{minipage} & \begin{minipage}[t]{0.05\columnwidth}\raggedright\strut
Re-exported exclusively from \texttt{runtime/app\_version.ts}; used for
version-gated migrations.\strut
\end{minipage}\tabularnewline
\begin{minipage}[t]{0.05\columnwidth}\raggedright\strut
\texttt{\_\_X2CHESS\_BUILD\_TIMESTAMP\_\_}\strut
\end{minipage} & \begin{minipage}[t]{0.05\columnwidth}\raggedright\strut
ISO-8601 string set at build time\strut
\end{minipage} & \begin{minipage}[t]{0.05\columnwidth}\raggedright\strut
Available for diagnostic display.\strut
\end{minipage}\tabularnewline
\bottomrule
\end{longtable}

All code that needs the version reads it from
\texttt{runtime/app\_version.ts}. No other file reads
\texttt{\_\_X2CHESS\_APP\_VERSION\_\_} directly.

\subsubsection{4.2 User preferences}\label{user-preferences}

User preferences are values the user can change within the application.
They persist across sessions via the versioned localStorage mechanism
described in §3.5.3.

Each preference group is a named versioned store backed by a single
localStorage key. Stores are read once on startup (after legacy key
consolidation; see §3.5.5) and written back whenever a preference
changes.

\paragraph{Preference stores}\label{preference-stores}

\begin{longtable}[]{@{}llll@{}}
\toprule
\begin{minipage}[b]{0.05\columnwidth}\raggedright\strut
Store\strut
\end{minipage} & \begin{minipage}[b]{0.05\columnwidth}\raggedright\strut
Key\strut
\end{minipage} & \begin{minipage}[b]{0.05\columnwidth}\raggedright\strut
File\strut
\end{minipage} & \begin{minipage}[b]{0.05\columnwidth}\raggedright\strut
Contents\strut
\end{minipage}\tabularnewline
\midrule
\endhead
\begin{minipage}[t]{0.05\columnwidth}\raggedright\strut
Shell prefs\strut
\end{minipage} & \begin{minipage}[t]{0.05\columnwidth}\raggedright\strut
\texttt{x2chess.shellPrefs}\strut
\end{minipage} & \begin{minipage}[t]{0.05\columnwidth}\raggedright\strut
\texttt{runtime/shell\_prefs\_store.ts}\strut
\end{minipage} & \begin{minipage}[t]{0.05\columnwidth}\raggedright\strut
Sound, locale, move delay, PGN layout mode, position preview, developer
tools, panel sizes\strut
\end{minipage}\tabularnewline
\begin{minipage}[t]{0.05\columnwidth}\raggedright\strut
Shape / board prefs\strut
\end{minipage} & \begin{minipage}[t]{0.05\columnwidth}\raggedright\strut
\texttt{x2chess.shapePrefs.v1}\strut
\end{minipage} & \begin{minipage}[t]{0.05\columnwidth}\raggedright\strut
\texttt{runtime/shape\_prefs.ts}\strut
\end{minipage} & \begin{minipage}[t]{0.05\columnwidth}\raggedright\strut
Annotation colors, square highlight style, move-hint dots\strut
\end{minipage}\tabularnewline
\begin{minipage}[t]{0.05\columnwidth}\raggedright\strut
External DB settings\strut
\end{minipage} & \begin{minipage}[t]{0.05\columnwidth}\raggedright\strut
\texttt{x2chess.ext-db-settings}\strut
\end{minipage} & \begin{minipage}[t]{0.05\columnwidth}\raggedright\strut
\texttt{runtime/ext\_database\_settings\_store.ts}\strut
\end{minipage} & \begin{minipage}[t]{0.05\columnwidth}\raggedright\strut
Opening explorer speed/rating filters\strut
\end{minipage}\tabularnewline
\begin{minipage}[t]{0.05\columnwidth}\raggedright\strut
Resource viewer column prefs\strut
\end{minipage} & \begin{minipage}[t]{0.05\columnwidth}\raggedright\strut
\texttt{x2chess.resourceViewerColumnPrefs.v1}\strut
\end{minipage} & \begin{minipage}[t]{0.05\columnwidth}\raggedright\strut
\texttt{resources\_viewer/viewer\_utils.ts}\strut
\end{minipage} & \begin{minipage}[t]{0.05\columnwidth}\raggedright\strut
Column widths and visibility\strut
\end{minipage}\tabularnewline
\bottomrule
\end{longtable}

For the complete list of all managed stores see the table in §3.5.3.

\paragraph{Preference lifecycle}\label{preference-lifecycle}

\begin{enumerate}
\def\labelenumi{\arabic{enumi}.}
\tightlist
\item
  \textbf{Startup} --- \texttt{useAppStartup} reads all stores
  synchronously; the values are dispatched to \texttt{AppStoreState} in
  the mount effect. Service-layer getters (\texttt{isSoundEnabled()},
  \texttt{getDelayMs()}, etc.) close over \texttt{stateRef}, so they
  pick up the current React state on every call without a separate sync
  step.
\item
  \textbf{Change} --- a user action dispatches through
  \texttt{ServiceContext}. The setter dispatches the new value to React
  and writes the store synchronously (\texttt{store.write(value)}).
  Service getters automatically see the new value on the next call via
  \texttt{stateRef}.
\item
  \textbf{Future build} --- if a newer build already wrote a
  higher-versioned envelope, the current build returns
  \texttt{defaultValue} without overwriting. The newer data is preserved
  for when that build runs again.
\end{enumerate}

\paragraph{Adding a preference}\label{adding-a-preference}

\begin{enumerate}
\def\labelenumi{\arabic{enumi}.}
\tightlist
\item
  Add the field to the relevant \texttt{*Prefs} type and to
  \texttt{DEFAULT\_*\_PREFS}.
\item
  If the stored shape changes incompatibly, increment \texttt{version}
  and append a \texttt{MigrationStep} to \texttt{migrations}.
\item
  Expose a setter callback in \texttt{ServiceContext} and wire it in
  \texttt{createAppServices.ts}.
\item
  Register the store in the table in §3.5.3 of this manual.
\end{enumerate}

\subsection{5) Resource Formats}\label{resource-formats}

X2Chess handles PGN content through three canonical resource kinds
(\texttt{file}, \texttt{directory}, \texttt{db}) plus a set of
import-time converters for external formats. All kinds share the same
adapter interface defined in \texttt{resource/domain/contracts.ts}; the
resource client in \texttt{resource/client/} routes every operation to
the correct adapter by \texttt{kind}.

The resource fomat \texttt{file} is a single PGN file with multiple
games. In this app this format is treated as read only as writing
requires a rewrite of the entire file. Writing ability may be
implemented in the future.

The resource fomat \texttt{directory} is a directory of PGN files. It is
implemented mainly for testing purposes, but it may also serve practical
use in situations where a user wishes to work with games in a storage
structure that is independent of database technology.

The resource fomat \texttt{db} is an X2Chess database. This type of
database is built using SQLite and is the primary editable resource
format. It is a self-contained file that contains all the data for a
single game library. It is the only format that supports the full range
of capabilities of the resource system, including search, filtering, and
sorting.

A resource is addressed by a \texttt{PgnResourceRef} (\texttt{kind} +
\texttt{locator}). A single game within a resource is addressed by a
\texttt{PgnGameRef} (\texttt{kind} + \texttt{locator} +
\texttt{recordId}).

\subsubsection{\texorpdfstring{5.1 PGN file (\texttt{file}
kind)}{5.1 PGN file (file kind)}}\label{pgn-file-file-kind}

\begin{itemize}
\tightlist
\item
  \textbf{Locator:} absolute filesystem path to a \texttt{.pgn} file.
\item
  \textbf{recordId:} 1-based integer position of the game within the
  file (as a string).
\item
  \textbf{Layout:} a single text file containing one or more PGN games
  concatenated in sequence. Games are split on
  \texttt{{[}Event\ "\ldots{}"{]}} header boundaries.
\item
  \textbf{Title derivation:} \texttt{Event:\ White\ -\ Black} →
  \texttt{White\ -\ Black} → \texttt{Event} → \texttt{Game\ N}.
\item
  \textbf{Capabilities:} \texttt{list} and \texttt{load} are
  implemented. \texttt{save} and \texttt{create} are deferred (the file
  format does not support partial writes without rewriting the whole
  file).
\item
  \textbf{Implementation:}
  \texttt{resource/adapters/file/file\_adapter.ts}.
\end{itemize}

\subsubsection{\texorpdfstring{5.2 PGN directory (\texttt{directory}
kind)}{5.2 PGN directory (directory kind)}}\label{pgn-directory-directory-kind}

\begin{itemize}
\tightlist
\item
  \textbf{Locator:} absolute filesystem path to a folder.
\item
  \textbf{recordId:} filename of the individual \texttt{.pgn} file
  within the folder (one game per file convention).
\item
  \textbf{Layout:} each file in the directory contains exactly one game.
  The display order and any extra metadata that PGN headers cannot carry
  are stored in a JSON sidecar file (\texttt{.x2chess-meta.json}) placed
  next to the directory. The sidecar schema is:
\end{itemize}

\texttt{json\ \ \ \{\ "version":\ 1,\ "games":\ \{\ "\textless{}filename\textgreater{}":\ \{\ "orderIndex":\ 0,\ "extra":\ \{\}\ \}\ \}\ \}}

The sidecar is read on every \texttt{list} call and written whenever
order or metadata changes. A missing or corrupt sidecar is silently
treated as empty. - \textbf{Capabilities:} \texttt{list}, \texttt{load},
\texttt{save}, \texttt{create}, and \texttt{reorder} are all delegated
to injected callbacks wired by \texttt{source\_gateway.ts}. -
\textbf{Implementation:}
\texttt{resource/adapters/directory/directory\_adapter.ts} +
\texttt{resource/adapters/directory/sidecar.ts}.

\subsubsection{\texorpdfstring{5.3 X2Chess database (\texttt{db}
kind)}{5.3 X2Chess database (db kind)}}\label{x2chess-database-db-kind}

The \texttt{db} kind is the primary editable resource format. It is a
SQLite file with a \texttt{.x2chess} extension. One file is one
self-contained game library; multiple libraries can be open
simultaneously (one \texttt{DbGateway} per file path).

\paragraph{5.3.1 File identity and
location}\label{file-identity-and-location}

\begin{itemize}
\tightlist
\item
  \textbf{Locator:} absolute filesystem path to a \texttt{.x2chess}
  SQLite file.
\item
  \textbf{recordId:} UUID string assigned at game creation time.
\item
  New databases are created as empty SQLite files; schema migrations run
  automatically on first access (idempotent runner in
  \texttt{resource/database/migrations/runner.ts}).
\end{itemize}

\paragraph{5.3.2 Schema}\label{schema}

The schema is managed by an ordered migration manifest
(\texttt{resource/database/schema/schema\_manifest.ts}). Each migration
is a version number plus a list of SQL statements executed in order.
Migrations are never modified after shipping; new steps are always
appended.

\textbf{\texttt{schema\_version}} (created by the migration runner
itself):

\begin{longtable}[]{@{}lll@{}}
\toprule
Column & Type & Description\tabularnewline
\midrule
\endhead
\texttt{version} & INTEGER PK & Applied migration version
number\tabularnewline
\texttt{applied\_at} & INTEGER & Unix timestamp (ms) when the migration
ran\tabularnewline
\bottomrule
\end{longtable}

\textbf{\texttt{games}} --- one row per stored game or position:

\begin{longtable}[]{@{}lll@{}}
\toprule
\begin{minipage}[b]{0.05\columnwidth}\raggedright\strut
Column\strut
\end{minipage} & \begin{minipage}[b]{0.05\columnwidth}\raggedright\strut
Type\strut
\end{minipage} & \begin{minipage}[b]{0.05\columnwidth}\raggedright\strut
Description\strut
\end{minipage}\tabularnewline
\midrule
\endhead
\begin{minipage}[t]{0.05\columnwidth}\raggedright\strut
\texttt{id}\strut
\end{minipage} & \begin{minipage}[t]{0.05\columnwidth}\raggedright\strut
TEXT PK\strut
\end{minipage} & \begin{minipage}[t]{0.05\columnwidth}\raggedright\strut
UUID assigned at creation\strut
\end{minipage}\tabularnewline
\begin{minipage}[t]{0.05\columnwidth}\raggedright\strut
\texttt{pgn\_text}\strut
\end{minipage} & \begin{minipage}[t]{0.05\columnwidth}\raggedright\strut
TEXT\strut
\end{minipage} & \begin{minipage}[t]{0.05\columnwidth}\raggedright\strut
Full PGN text including headers and move tree\strut
\end{minipage}\tabularnewline
\begin{minipage}[t]{0.05\columnwidth}\raggedright\strut
\texttt{title\_hint}\strut
\end{minipage} & \begin{minipage}[t]{0.05\columnwidth}\raggedright\strut
TEXT\strut
\end{minipage} & \begin{minipage}[t]{0.05\columnwidth}\raggedright\strut
Fallback display title when PGN headers are absent\strut
\end{minipage}\tabularnewline
\begin{minipage}[t]{0.05\columnwidth}\raggedright\strut
\texttt{created\_at}\strut
\end{minipage} & \begin{minipage}[t]{0.05\columnwidth}\raggedright\strut
INTEGER\strut
\end{minipage} & \begin{minipage}[t]{0.05\columnwidth}\raggedright\strut
Unix timestamp (ms) of first insert\strut
\end{minipage}\tabularnewline
\begin{minipage}[t]{0.05\columnwidth}\raggedright\strut
\texttt{updated\_at}\strut
\end{minipage} & \begin{minipage}[t]{0.05\columnwidth}\raggedright\strut
INTEGER\strut
\end{minipage} & \begin{minipage}[t]{0.05\columnwidth}\raggedright\strut
Unix timestamp (ms) of last save\strut
\end{minipage}\tabularnewline
\begin{minipage}[t]{0.05\columnwidth}\raggedright\strut
\texttt{order\_index}\strut
\end{minipage} & \begin{minipage}[t]{0.05\columnwidth}\raggedright\strut
REAL\strut
\end{minipage} & \begin{minipage}[t]{0.05\columnwidth}\raggedright\strut
Display order; swap-reorder exchanges two values\strut
\end{minipage}\tabularnewline
\begin{minipage}[t]{0.05\columnwidth}\raggedright\strut
\texttt{revision\_token}\strut
\end{minipage} & \begin{minipage}[t]{0.05\columnwidth}\raggedright\strut
TEXT\strut
\end{minipage} & \begin{minipage}[t]{0.05\columnwidth}\raggedright\strut
Opaque token updated on every save for conflict detection\strut
\end{minipage}\tabularnewline
\begin{minipage}[t]{0.05\columnwidth}\raggedright\strut
\texttt{kind}\strut
\end{minipage} & \begin{minipage}[t]{0.05\columnwidth}\raggedright\strut
TEXT\strut
\end{minipage} & \begin{minipage}[t]{0.05\columnwidth}\raggedright\strut
\texttt{"game"} (default) or \texttt{"position"} (detected from
\texttt{{[}SetUp\ "1"{]}})\strut
\end{minipage}\tabularnewline
\bottomrule
\end{longtable}

\textbf{\texttt{game\_metadata}} --- denormalized header values for fast
search and display:

\begin{longtable}[]{@{}lll@{}}
\toprule
Column & Type & Description\tabularnewline
\midrule
\endhead
\texttt{game\_id} & TEXT FK → games.id & Parent game (cascades on
delete)\tabularnewline
\texttt{meta\_key} & TEXT & PGN header key (e.g. \texttt{White},
\texttt{Black}, \texttt{Event}, \texttt{Result})\tabularnewline
\texttt{val\_str} & TEXT & String value extracted from the PGN
header\tabularnewline
\bottomrule
\end{longtable}

Indexed on \texttt{(meta\_key)} and \texttt{(meta\_key,\ val\_str)} for
text-search queries.

\textbf{\texttt{metadata\_keys}} --- registry of user-defined metadata
key names:

\begin{longtable}[]{@{}lll@{}}
\toprule
Column & Type & Description\tabularnewline
\midrule
\endhead
\texttt{key} & TEXT PK & Metadata key name\tabularnewline
\texttt{value\_type} & TEXT & Value type hint (\texttt{"string"}
default)\tabularnewline
\bottomrule
\end{longtable}

\textbf{\texttt{position\_hashes}} --- position search index (one row
per reachable position per game):

\begin{longtable}[]{@{}lll@{}}
\toprule
Column & Type & Description\tabularnewline
\midrule
\endhead
\texttt{hash} & TEXT & 16-char hex FNV-1a hash of the normalized FEN
(first 4 fields)\tabularnewline
\texttt{game\_id} & TEXT FK → games.id & Parent game (cascades on
delete)\tabularnewline
\texttt{ply} & INTEGER & Half-move number where this position occurs (0
= start)\tabularnewline
\bottomrule
\end{longtable}

Indexed on \texttt{hash} for position-lookup queries. The hash function
strips half-move and full-move counters so transpositions hash
identically. The index is rebuilt in full on every save; the builder is
injectable (\texttt{BuildPositionIndex}) so the \texttt{resource/}
package stays free of chess-engine dependencies.

\textbf{\texttt{move\_edges}} --- move frequency index for opening
explorer:

\begin{longtable}[]{@{}lll@{}}
\toprule
Column & Type & Description\tabularnewline
\midrule
\endhead
\texttt{position\_hash} & TEXT & Position hash (same scheme as
\texttt{position\_hashes.hash})\tabularnewline
\texttt{move\_san} & TEXT & Move in SAN notation\tabularnewline
\texttt{move\_uci} & TEXT & Move in UCI notation\tabularnewline
\texttt{result} & TEXT & Game result (\texttt{1-0}, \texttt{0-1},
\texttt{1/2-1/2}, \texttt{*})\tabularnewline
\texttt{game\_id} & TEXT FK → games.id & Parent game (cascades on
delete)\tabularnewline
\bottomrule
\end{longtable}

Indexed on \texttt{position\_hash}. Used by \texttt{explorePosition} to
return per-move frequency/result aggregates for any reachable position
across the whole library.

\textbf{\texttt{training\_transcripts}} --- per-training-session records
(owned by the training module):

\begin{longtable}[]{@{}lll@{}}
\toprule
Column & Type & Description\tabularnewline
\midrule
\endhead
\texttt{session\_id} & TEXT PK & UUID for the training
session\tabularnewline
\texttt{source\_game\_id} & TEXT & The game being trained
against\tabularnewline
\texttt{protocol} & TEXT & Training protocol identifier\tabularnewline
\texttt{started\_at} & INTEGER & Unix timestamp (ms)\tabularnewline
\texttt{completed\_at} & INTEGER & Unix timestamp (ms),
nullable\tabularnewline
\texttt{aborted} & INTEGER & 0 = not aborted, 1 = aborted\tabularnewline
\texttt{config\_json} & TEXT & JSON-serialized training
configuration\tabularnewline
\texttt{ply\_records} & TEXT & JSON array of per-ply training
results\tabularnewline
\texttt{annotations} & TEXT & JSON array of post-session
annotations\tabularnewline
\bottomrule
\end{longtable}

Indexed on \texttt{source\_game\_id} for per-game transcript queries.

\paragraph{5.3.3 Write path and revision
safety}\label{write-path-and-revision-safety}

Every \texttt{save} call: 1. Optionally checks \texttt{revision\_token}
against the stored value; throws \texttt{conflict} on mismatch
(optimistic concurrency guard). 2. Updates \texttt{pgn\_text},
\texttt{revision\_token}, \texttt{updated\_at}, and \texttt{kind} in
\texttt{games}. 3. Deletes and rebuilds \texttt{game\_metadata},
\texttt{position\_hashes}, and \texttt{move\_edges} for the affected
game.

Every \texttt{create} call assigns a UUID, timestamps the row, appends
\texttt{order\_index} after the current maximum, and builds all three
indexes.

\paragraph{5.3.4 Capabilities}\label{capabilities}

\begin{longtable}[]{@{}ll@{}}
\toprule
Operation & Supported\tabularnewline
\midrule
\endhead
\texttt{list} & Yes --- ordered by
\texttt{order\_index\ ASC,\ created\_at\ ASC}\tabularnewline
\texttt{load} & Yes\tabularnewline
\texttt{save} & Yes --- with optional revision-token conflict
detection\tabularnewline
\texttt{create} & Yes\tabularnewline
\texttt{reorder} & Yes --- swaps \texttt{order\_index} values between
two rows\tabularnewline
\texttt{searchByPositionHash} & Yes --- queries
\texttt{position\_hashes}\tabularnewline
\texttt{searchByText} & Yes --- substring match on \texttt{White},
\texttt{Black}, \texttt{Event}, \texttt{Site}\tabularnewline
\texttt{explorePosition} & Yes --- aggregated move frequency from
\texttt{move\_edges}\tabularnewline
\bottomrule
\end{longtable}

\subsubsection{5.4 Import-time format
converters}\label{import-time-format-converters}

Import converters are one-shot adapters that translate a foreign file
format into PGN strings. They implement \texttt{FormatImporter}
(\texttt{resource/adapters/import/format\_import\_types.ts}) and are not
resource adapters --- they do not implement
\texttt{list}/\texttt{load}/\texttt{save}. Converted games are handed to
the caller as PGN text ready to be stored in a \texttt{db} or other
resource.

\begin{longtable}[]{@{}lll@{}}
\toprule
\begin{minipage}[b]{0.05\columnwidth}\raggedright\strut
Format\strut
\end{minipage} & \begin{minipage}[b]{0.05\columnwidth}\raggedright\strut
Extensions\strut
\end{minipage} & \begin{minipage}[b]{0.05\columnwidth}\raggedright\strut
Mechanism\strut
\end{minipage}\tabularnewline
\midrule
\endhead
\begin{minipage}[t]{0.05\columnwidth}\raggedright\strut
EPD (Extended Position Description)\strut
\end{minipage} & \begin{minipage}[t]{0.05\columnwidth}\raggedright\strut
\texttt{.epd}\strut
\end{minipage} & \begin{minipage}[t]{0.05\columnwidth}\raggedright\strut
Pure TypeScript parser; each line becomes a position-only PGN game with
\texttt{{[}SetUp\ "1"{]}} and \texttt{{[}FEN\ "\ldots{}"{]}} headers.
Opcodes \texttt{id}, \texttt{bm}, \texttt{c0}--\texttt{c9}, and
\texttt{ce} are mapped to PGN headers and comments.\strut
\end{minipage}\tabularnewline
\begin{minipage}[t]{0.05\columnwidth}\raggedright\strut
ChessBase\strut
\end{minipage} & \begin{minipage}[t]{0.05\columnwidth}\raggedright\strut
\texttt{.cbh}, \texttt{.cbv}\strut
\end{minipage} & \begin{minipage}[t]{0.05\columnwidth}\raggedright\strut
Delegates binary parsing to the Tauri backend command
\texttt{import\_chessbase\_file}; Rust returns an array of PGN strings.
The TypeScript layer is complete; the Rust command is deferred.\strut
\end{minipage}\tabularnewline
\bottomrule
\end{longtable}

\textbf{EPD opcode mapping:}

\begin{longtable}[]{@{}ll@{}}
\toprule
Opcode & PGN mapping\tabularnewline
\midrule
\endhead
\texttt{id} & \texttt{{[}Event\ "\ldots{}"{]}} header\tabularnewline
\texttt{bm} & \texttt{{[}Annotator\ "bm:\ \ldots{}"{]}}
header\tabularnewline
\texttt{c0}--\texttt{c9} & Comment block \texttt{\{\ \ldots{}\ \}} in
move text\tabularnewline
\texttt{ce} & Centipawn eval appended to comment:
\texttt{eval:\ \textless{}value\textgreater{}}\tabularnewline
All others & Ignored\tabularnewline
\bottomrule
\end{longtable}

\subsection{6) Testing}\label{testing}

\subsubsection{5.1 Test runner and
location}\label{test-runner-and-location}

Tests use Node's built-in test runner (\texttt{node:test}) invoked
through \texttt{tsx} for TypeScript support. No third-party test
framework or mocking library is used; tests rely on real implementations
or hand-written stubs.

\begin{Shaded}
\begin{Highlighting}[]
\CommentTok{# from frontend/}
\KeywordTok{npm} \NormalTok{test                                       }\CommentTok{# run all tests}
\KeywordTok{tsx} \NormalTok{--test }\StringTok{"test/foo/**/*.test.ts"}             \CommentTok{# run a subset}
\end{Highlighting}
\end{Shaded}

Test files live under \texttt{frontend/test/}, mirroring the
\texttt{src/} structure exactly --- for example
\texttt{src/board/shape\_serializer.ts} →
\texttt{test/board/shape\_serializer.test.ts}.

\subsubsection{5.2 What to test}\label{what-to-test}

\textbf{Pure-logic modules are the primary test target.} Any module
under \texttt{model/}, \texttt{editor/}, \texttt{board/},
\texttt{game\_sessions/}, \texttt{resources/},
\texttt{resources\_viewer/}, \texttt{runtime/}, or \texttt{training/}
that exports non-trivial logic should have a corresponding test file.

The \texttt{state/app\_reducer.ts} reducer and
\texttt{state/selectors.ts} are pure functions and are tested like any
other pure-logic module, despite living in \texttt{state/}.

\textbf{React components, hooks, and Tauri gateways are not
unit-tested.} They belong to the integration layer and are validated
through manual QA.

\subsubsection{5.3 Coverage map}\label{coverage-map}

\begin{longtable}[]{@{}ll@{}}
\toprule
\begin{minipage}[b]{0.05\columnwidth}\raggedright\strut
Test directory\strut
\end{minipage} & \begin{minipage}[b]{0.05\columnwidth}\raggedright\strut
Source coverage\strut
\end{minipage}\tabularnewline
\midrule
\endhead
\begin{minipage}[t]{0.05\columnwidth}\raggedright\strut
\texttt{test/board/}\strut
\end{minipage} & \begin{minipage}[t]{0.05\columnwidth}\raggedright\strut
\texttt{src/board/} --- move lookup, navigation, shape
parser/serializer, physical board protocols\strut
\end{minipage}\tabularnewline
\begin{minipage}[t]{0.05\columnwidth}\raggedright\strut
\texttt{test/editor/}\strut
\end{minipage} & \begin{minipage}[t]{0.05\columnwidth}\raggedright\strut
\texttt{src/editor/} --- FEN utils, tree numbering, collapse, editor
plan, anchor resolver\strut
\end{minipage}\tabularnewline
\begin{minipage}[t]{0.05\columnwidth}\raggedright\strut
\texttt{test/engines/}\strut
\end{minipage} & \begin{minipage}[t]{0.05\columnwidth}\raggedright\strut
Engine manager, UCI parser/writer/session\strut
\end{minipage}\tabularnewline
\begin{minipage}[t]{0.05\columnwidth}\raggedright\strut
\texttt{test/game\_sessions/}\strut
\end{minipage} & \begin{minipage}[t]{0.05\columnwidth}\raggedright\strut
Session store, persistence, ingress handlers\strut
\end{minipage}\tabularnewline
\begin{minipage}[t]{0.05\columnwidth}\raggedright\strut
\texttt{test/model/}\strut
\end{minipage} & \begin{minipage}[t]{0.05\columnwidth}\raggedright\strut
PGN model parser, serializer, commands, move ops, undo/redo,
x2-style\strut
\end{minipage}\tabularnewline
\begin{minipage}[t]{0.05\columnwidth}\raggedright\strut
\texttt{test/resources/}\strut
\end{minipage} & \begin{minipage}[t]{0.05\columnwidth}\raggedright\strut
Material key, PGN DB adapter, PGN metadata, web-import rules\strut
\end{minipage}\tabularnewline
\begin{minipage}[t]{0.05\columnwidth}\raggedright\strut
\texttt{test/resources\_viewer/}\strut
\end{minipage} & \begin{minipage}[t]{0.05\columnwidth}\raggedright\strut
Annotation parsers: anchor, eval, link, QA, schema storage, todo, train
tag\strut
\end{minipage}\tabularnewline
\begin{minipage}[t]{0.05\columnwidth}\raggedright\strut
\texttt{test/runtime/}\strut
\end{minipage} & \begin{minipage}[t]{0.05\columnwidth}\raggedright\strut
Resource ref utilities\strut
\end{minipage}\tabularnewline
\begin{minipage}[t]{0.05\columnwidth}\raggedright\strut
\texttt{test/state/}\strut
\end{minipage} & \begin{minipage}[t]{0.05\columnwidth}\raggedright\strut
App reducer, selectors\strut
\end{minipage}\tabularnewline
\begin{minipage}[t]{0.05\columnwidth}\raggedright\strut
\texttt{test/training/}\strut
\end{minipage} & \begin{minipage}[t]{0.05\columnwidth}\raggedright\strut
Move acceptance, replay protocol, transcript, merge, train-tag parser,
curriculum I/O\strut
\end{minipage}\tabularnewline
\bottomrule
\end{longtable}

\subsubsection{5.4 Writing new tests}\label{writing-new-tests}

\begin{itemize}
\tightlist
\item
  Import source files with the \texttt{.js} extension (required by the
  ESM/tsx loader):
  \texttt{import\ \{\ foo\ \}\ from\ "../../src/bar.js"}.
\item
  Use \texttt{String.raw\textbackslash{}}\ldots{}`` for strings
  containing backslash sequences (PGN escape characters, regex literals)
  to avoid double-escaping noise.
\item
  Stub external dependencies with minimal inline objects; never
  duplicate production logic in stubs.
\item
  One test file per source module --- do not mix multiple modules into
  one file.
\end{itemize}

\subsection{7) Key Concepts}\label{key-concepts}

\subsubsection{7.1 GameSessionState and
ActiveSessionRef}\label{gamesessionstate-and-activesessionref}

Each open game tab is backed by a dedicated \texttt{GameSessionState}
object (\texttt{frontend/src/game\_sessions/game\_session\_state.ts}).
This object is the \textbf{working store} for all game-specific mutable
state for one session:

\begin{itemize}
\tightlist
\item
  \textbf{PGN / model}: \texttt{pgnModel}, \texttt{pgnText},
  \texttt{moves}, \texttt{verboseMoves}, \texttt{movePositionById},
  \texttt{pgnLayoutMode}
\item
  \textbf{Navigation}: \texttt{currentPly}, \texttt{selectedMoveId},
  \texttt{boardPreview}
\item
  \textbf{Editor messages}: \texttt{pendingFocusCommentId}
\item
  \textbf{Edit history}: \texttt{undoStack}, \texttt{redoStack}
\end{itemize}

Fields that are internal to a specific service and never need to cross
service boundaries live as closure-private variables instead of being
stored on \texttt{GameSessionState}:

\begin{itemize}
\tightlist
\item
  \texttt{animationRunId}, \texttt{isAnimating} --- owned by the
  \texttt{navigation} closure
\item
  \texttt{errorMessage}, \texttt{statusMessage} --- dispatched directly
  as \texttt{set\_error\_message} / \texttt{set\_status\_message}
  actions; not stored on the session object
\end{itemize}

Service modules (\texttt{pgn\_runtime}, \texttt{history},
\texttt{move\_lookup}, \texttt{navigation}, \texttt{pgn\_model\_update})
never receive a global state object. They receive an
\texttt{ActiveSessionRef\ =\ \{\ current:\ GameSessionState\ \}} and
read/write only \texttt{sessionRef.current.*}. This makes the active
session swappable without copying data: a session switch is a single
pointer assignment
\texttt{activeSessionRef.current\ =\ target.ownState}. Each
\texttt{GameSession} in the session store holds its live
\texttt{GameSessionState} in an \texttt{ownState} field --- there is no
freeze/thaw snapshot cycle.

\subsubsection{7.2 Service-Owned State}\label{service-owned-state}

Each service manages its own internal state as closure-private
variables. There is no shared mutable \texttt{AppState} object.

\begin{longtable}[]{@{}ll@{}}
\toprule
\begin{minipage}[b]{0.05\columnwidth}\raggedright\strut
Service\strut
\end{minipage} & \begin{minipage}[b]{0.05\columnwidth}\raggedright\strut
Owns\strut
\end{minipage}\tabularnewline
\midrule
\endhead
\begin{minipage}[t]{0.05\columnwidth}\raggedright\strut
\texttt{session\_store}\strut
\end{minipage} & \begin{minipage}[t]{0.05\columnwidth}\raggedright\strut
\texttt{gameSessions{[}{]}}, \texttt{activeSessionId},
\texttt{nextSessionSeq}\strut
\end{minipage}\tabularnewline
\begin{minipage}[t]{0.05\columnwidth}\raggedright\strut
\texttt{resources\_viewer}\strut
\end{minipage} & \begin{minipage}[t]{0.05\columnwidth}\raggedright\strut
\texttt{resourceViewerTabs{[}{]}}, \texttt{activeResourceTabId},
metadata-key prefs\strut
\end{minipage}\tabularnewline
\begin{minipage}[t]{0.05\columnwidth}\raggedright\strut
\texttt{session\_persistence}\strut
\end{minipage} & \begin{minipage}[t]{0.05\columnwidth}\raggedright\strut
\texttt{autosaveTimer}, \texttt{saveRequestSeq},
\texttt{isHydratingGame}, \texttt{defaultSaveMode}\strut
\end{minipage}\tabularnewline
\begin{minipage}[t]{0.05\columnwidth}\raggedright\strut
\texttt{resources}\strut
\end{minipage} & \begin{minipage}[t]{0.05\columnwidth}\raggedright\strut
\texttt{gameDirectoryHandle}, paths, \texttt{activeSourceKind},
\texttt{appConfig}, \texttt{playerStore}\strut
\end{minipage}\tabularnewline
\begin{minipage}[t]{0.05\columnwidth}\raggedright\strut
\texttt{navigation}\strut
\end{minipage} & \begin{minipage}[t]{0.05\columnwidth}\raggedright\strut
\texttt{animationRunId}, \texttt{isAnimating}\strut
\end{minipage}\tabularnewline
\begin{minipage}[t]{0.05\columnwidth}\raggedright\strut
\texttt{history}\strut
\end{minipage} & \begin{minipage}[t]{0.05\columnwidth}\raggedright\strut
(undo/redo stacks are on \texttt{GameSessionState} so they survive
session switches)\strut
\end{minipage}\tabularnewline
\bottomrule
\end{longtable}

Services that need to read shell-preference values (sound, locale, move
delay) receive a getter function injected at bundle-creation time:
\texttt{isSoundEnabled()}, \texttt{getDelayMs()}, etc. These getters
close over a \texttt{stateRef} maintained in \texttt{useAppStartup} that
is updated on every React render, so services always read the current
preference without coupling to React directly.

Shell preferences (sound, locale, move delay, layout mode, position
preview, developer tools) are the authoritative province of the React
store (\texttt{AppStoreState}) and are persisted via
\texttt{shellPrefsStore} --- a versioned localStorage store
(\texttt{frontend/src/runtime/shell\_prefs\_store.ts}). Shell preference
setters dispatch to React directly.

\subsubsection{7.3 The React State Bridge}\label{the-react-state-bridge}

The React \texttt{useReducer} store (\texttt{AppStoreState} /
\texttt{appReducer}) is the authoritative state for \textbf{rendering}.
State reaches it through two channels:

\paragraph{Service callbacks
(automatic)}\label{service-callbacks-automatic}

Services wired in \texttt{createAppServices.ts} call typed dispatch
callbacks when their state changes:

\begin{longtable}[]{@{}ll@{}}
\toprule
Callback & Action dispatched\tabularnewline
\midrule
\endhead
\texttt{onPgnChange} & \texttt{set\_pgn\_state}\tabularnewline
\texttt{onNavigationChange} & \texttt{set\_navigation}\tabularnewline
\texttt{onUndoRedoDepthChange} &
\texttt{set\_undo\_redo\_depth}\tabularnewline
\texttt{onStatusChange} & \texttt{set\_status\_message}\tabularnewline
\texttt{onErrorChange} & \texttt{set\_error\_message}\tabularnewline
\texttt{onPendingFocusChange} &
\texttt{set\_pending\_focus}\tabularnewline
\texttt{onSessionsChanged} & \texttt{set\_sessions}\tabularnewline
\texttt{onTabsChanged} & \texttt{set\_resource\_viewer}\tabularnewline
\bottomrule
\end{longtable}

Callbacks close over \texttt{dispatchRef} (a mutable ref updated on
every render) so they always use the current React dispatch without
recreating service instances.

\paragraph{flushSessionState (explicit, for direct
mutations)}\label{flushsessionstate-explicit-for-direct-mutations}

Operations in \texttt{SessionOrchestrator} that mutate
\texttt{activeSessionRef.current} directly (e.g. \texttt{loadPgnText},
\texttt{switchSession}, \texttt{openPgnText}) call
\texttt{flushSessionState()} afterwards. This helper reads the active
session and dispatches \texttt{set\_pgn\_state},
\texttt{set\_navigation}, \texttt{set\_undo\_redo\_depth}, and
\texttt{set\_pending\_focus} in one call.

\paragraph{The mutable working store}\label{the-mutable-working-store}

\texttt{activeSessionRef.current} (\texttt{GameSessionState}) is a
\textbf{deliberate mutable working store} --- not a React anti-pattern.
Operations mutate it in-place during a single logical step (e.g.~parse
PGN, update ply, update selection), then flush to React once the step is
complete. This avoids mid-operation renders and keeps animation, undo,
and board preview coherent.

The pattern is intentional: - Components \textbf{read} from the React
store (always consistent). - Services \textbf{write} to
\texttt{activeSessionRef.current} and flush when done. - The React store
is the display layer; \texttt{activeSessionRef.current} is the editing
scratch pad for the active session.

\paragraph{Component contract}\label{component-contract}

React components read state exclusively through selectors
(\texttt{state/selectors.ts}). Components trigger mutations exclusively
via \texttt{ServiceContext} callbacks --- they never call
\texttt{dispatch} for game logic.

\subsubsection{7.4 The PGN Text Editor}\label{the-pgn-text-editor}

The PGN text editor is the annotation-aware view of the active game
session's PGN model. It is structured as a pure-logic
\textbf{render-plan pipeline} that converts the in-memory PGN model into
a flat list of \texttt{PlanBlock{[}{]}}, which the React render layer
then maps to DOM elements.

\paragraph{Layout modes}\label{layout-modes}

The editor operates in one of three layout modes, stored per session in
\texttt{GameSessionState.pgnLayoutMode}:

\begin{longtable}[]{@{}ll@{}}
\toprule
\begin{minipage}[b]{0.05\columnwidth}\raggedright\strut
Mode\strut
\end{minipage} & \begin{minipage}[b]{0.05\columnwidth}\raggedright\strut
Behaviour\strut
\end{minipage}\tabularnewline
\midrule
\endhead
\begin{minipage}[t]{0.05\columnwidth}\raggedright\strut
\texttt{plain}\strut
\end{minipage} & \begin{minipage}[t]{0.05\columnwidth}\raggedright\strut
Raw PGN text; comment content is shown verbatim, no marker
processing.\strut
\end{minipage}\tabularnewline
\begin{minipage}[t]{0.05\columnwidth}\raggedright\strut
\texttt{text}\strut
\end{minipage} & \begin{minipage}[t]{0.05\columnwidth}\raggedright\strut
\texttt{{[}{[}indent{]}{]}} drives block indentation;
\texttt{{[}{[}br{]}{]}} becomes a visible line break in the
\texttt{contentEditable}; the first comment before the first move
receives \emph{intro styling}.\strut
\end{minipage}\tabularnewline
\begin{minipage}[t]{0.05\columnwidth}\raggedright\strut
\texttt{tree}\strut
\end{minipage} & \begin{minipage}[t]{0.05\columnwidth}\raggedright\strut
One block per variation in DFS order; each non-mainline block opens with
a labelled \texttt{branch\_header} token; markers survive as greyed
literal text; each variation's first comment receives intro styling
independently.\strut
\end{minipage}\tabularnewline
\bottomrule
\end{longtable}

\paragraph{Intro section (implicit
detection)}\label{intro-section-implicit-detection}

When the editor is in \textbf{text} or \textbf{tree} mode, the
\emph{first comment appearing before the first move of a variation}
automatically receives \textbf{intro styling} --- a visually distinct
block with a configurable left sidebar and tinted background. No PGN
marker is required: position in the model is the sole criterion.

Detection inside \texttt{buildVariationWalker} (\texttt{plan/types.ts}):

\begin{verbatim}
const applyIntroStyling: boolean = isFirstComment && !state.firstMoveEmitted;
\end{verbatim}

\begin{itemize}
\tightlist
\item
  \texttt{firstCommentId} --- set on the first comment seen in the
  current scope; a later comment in the same variation cannot be the
  intro.
\item
  \texttt{firstMoveEmitted} --- set to \texttt{true} once any move token
  has been emitted; a comment appearing after the first move is never
  styled as intro.
\end{itemize}

In \textbf{tree mode} \texttt{buildTreeEditorPlan} saves and resets both
fields before processing each variation block, so \textbf{each
variation's first comment gets intro styling independently} of the
main-line intro.

In \textbf{plain mode} \texttt{applyIntroStyling} is always
\texttt{false}; the detection is skipped entirely and all comments
render as plain literal text.

\textbf{CSS}: the class \texttt{text-editor-comment-intro} is added when
\texttt{CommentToken.introStyling\ ===\ true}. Visual appearance is
driven by the CSS custom properties \texttt{-\/-editor-intro-bg},
\texttt{-\/-editor-intro-sidebar-color}, and
\texttt{-\/-editor-intro-sidebar-width}, which are set from
\texttt{EditorStylePrefs} and default to the shipped design.

\begin{quote}
\textbf{Note:} The legacy \texttt{\textbackslash{}intro} PGN marker is
gone. Intro styling is fully implicit --- no markup is needed in the PGN
source.
\end{quote}

\paragraph{\texorpdfstring{Render-plan pipeline
(\texttt{editor/plan/})}{Render-plan pipeline (editor/plan/)}}\label{render-plan-pipeline-editorplan}

\texttt{buildTextEditorPlan(pgnModel,\ \{\ layoutMode,\ numberingStrategy\ \})}
is the single entry point (re-exported from
\texttt{editor/text\_editor\_plan.ts}). It:

\begin{enumerate}
\def\labelenumi{\arabic{enumi}.}
\tightlist
\item
  Allocates a \texttt{PlanState} --- a cursor over a growing
  \texttt{PlanBlock{[}{]}} plus tracking fields (\texttt{indentDepth},
  \texttt{firstCommentId}, \texttt{firstMoveEmitted}).
\item
  Delegates to the mode-specific builder:
\end{enumerate}

\begin{itemize}
\tightlist
\item
  \texttt{buildPlainEditorPlan} (\texttt{plan/plain\_mode.ts}) --- uses
  \texttt{buildVariationWalker} with a pass-through comment emitter that
  sets \texttt{plainLiteralComment\ =\ true}.
\item
  \texttt{buildTextModeEditorPlan} (\texttt{plan/text\_mode.ts}) ---
  uses \texttt{buildVariationWalker} with a comment emitter that strips
  \texttt{{[}{[}indent{]}{]}} directives, converts
  \texttt{{[}{[}br{]}{]}} to \texttt{\textbackslash{}n}, and calls
  \texttt{nextBlock()} after intro comments.
\item
  \texttt{buildTreeEditorPlan} (\texttt{plan/tree\_mode.ts}) --- has its
  own DFS traversal; collects child RAVs inline and emits them as new
  blocks after the parent block's content.
\end{itemize}

\begin{enumerate}
\def\labelenumi{\arabic{enumi}.}
\setcounter{enumi}{2}
\tightlist
\item
  Trims leading and trailing empty blocks before returning.
\end{enumerate}

\textbf{\texttt{buildVariationWalker}} (\texttt{plan/types.ts}) --- a
factory used by plain and text mode. It returns a mutually-recursive
\texttt{emitVariation} / \texttt{emitMove} pair bound to a mode-specific
\texttt{CommentEmitter}. The walker handles comment hoisting (a comment
immediately before a move number is hoisted before that move's token),
\texttt{{[}{[}indent{]}{]}}- gated RAV indentation, and
\texttt{postItems} ordering from the parsed PGN model.

\textbf{Token types} emitted into blocks:

\begin{longtable}[]{@{}ll@{}}
\toprule
\begin{minipage}[b]{0.05\columnwidth}\raggedright\strut
\texttt{tokenType}\strut
\end{minipage} & \begin{minipage}[b]{0.05\columnwidth}\raggedright\strut
Description\strut
\end{minipage}\tabularnewline
\midrule
\endhead
\begin{minipage}[t]{0.05\columnwidth}\raggedright\strut
\texttt{move}\strut
\end{minipage} & \begin{minipage}[t]{0.05\columnwidth}\raggedright\strut
A SAN move token; carries \texttt{nodeId}, \texttt{variationDepth},
\texttt{moveSide}.\strut
\end{minipage}\tabularnewline
\begin{minipage}[t]{0.05\columnwidth}\raggedright\strut
\texttt{move\_number}\strut
\end{minipage} & \begin{minipage}[t]{0.05\columnwidth}\raggedright\strut
A move-number token (simplified to digit only).\strut
\end{minipage}\tabularnewline
\begin{minipage}[t]{0.05\columnwidth}\raggedright\strut
\texttt{nag}\strut
\end{minipage} & \begin{minipage}[t]{0.05\columnwidth}\raggedright\strut
A NAG glyph (e.g. \texttt{!}, \texttt{?}, \texttt{±}).\strut
\end{minipage}\tabularnewline
\begin{minipage}[t]{0.05\columnwidth}\raggedright\strut
\texttt{result}\strut
\end{minipage} & \begin{minipage}[t]{0.05\columnwidth}\raggedright\strut
The game result (\texttt{1-0}, \texttt{0-1}, \texttt{½-½},
\texttt{*}).\strut
\end{minipage}\tabularnewline
\begin{minipage}[t]{0.05\columnwidth}\raggedright\strut
\texttt{space}\strut
\end{minipage} & \begin{minipage}[t]{0.05\columnwidth}\raggedright\strut
A whitespace spacer between tokens.\strut
\end{minipage}\tabularnewline
\begin{minipage}[t]{0.05\columnwidth}\raggedright\strut
\texttt{branch\_header}\strut
\end{minipage} & \begin{minipage}[t]{0.05\columnwidth}\raggedright\strut
Tree-mode only; carries the variation label (e.g. \texttt{A},
\texttt{B.1}).\strut
\end{minipage}\tabularnewline
\begin{minipage}[t]{0.05\columnwidth}\raggedright\strut
\texttt{comment} (kind)\strut
\end{minipage} & \begin{minipage}[t]{0.05\columnwidth}\raggedright\strut
A \texttt{CommentToken}; carries \texttt{commentId}, \texttt{rawText},
display \texttt{text}, indent flags, \texttt{introStyling}, and
\texttt{focusFirstCommentAtStart}.\strut
\end{minipage}\tabularnewline
\bottomrule
\end{longtable}

\textbf{Tree variation numbering} (\texttt{plan/tree\_mode.ts}): each
non-mainline block is labelled by a \texttt{VariationNumberingStrategy}
--- a function from a path array to a string. The default
\texttt{alphaNumericPathStrategy} maps the first segment to an uppercase
letter (\texttt{A}, \texttt{B}, \ldots{}) and deeper segments to 1-based
integers joined by \texttt{.} (e.g. \texttt{A.1.2}).

\paragraph{\texorpdfstring{PGN runtime
(\texttt{editor/pgn\_runtime.ts})}{PGN runtime (editor/pgn\_runtime.ts)}}\label{pgn-runtime-editorpgn_runtime.ts}

\texttt{createPgnRuntimeCapabilities} owns the lifecycle of the
in-session PGN text:

\begin{itemize}
\tightlist
\item
  \textbf{\texttt{loadPgn()}} --- reads the raw PGN text from the editor
  input element, parses it into a model, runs
  \texttt{syncChessParseState}, and calls \texttt{onPgnChange} +
  \texttt{onNavigationChange}.
\item
  \textbf{\texttt{applyPgnModelUpdate(nextModel,\ focusCommentId?)}} ---
  the primary mutation path for model-level edits (commands, move
  entry). Serializes the new model to text, pushes an undo snapshot if
  the text changed, syncs chess parse state, schedules autosave, and
  calls \texttt{onPgnChange} + \texttt{onNavigationChange} + (if
  applicable) \texttt{onPendingFocusChange}.
\item
  \textbf{\texttt{syncChessParseState(source)}} --- feeds the PGN text
  to \texttt{chess.js} to populate \texttt{moves},
  \texttt{verboseMoves}, and \texttt{movePositionById}. Falls back to a
  stripped (annotation-free) copy when annotations cause
  \texttt{chess.js} to reject the input. Does not fire any callbacks;
  callers dispatch state separately.
\end{itemize}

\paragraph{\texorpdfstring{Undo / redo
(\texttt{editor/history.ts})}{Undo / redo (editor/history.ts)}}\label{undo-redo-editorhistory.ts}

\texttt{createEditorHistoryCapabilities} maintains per-session
\texttt{undoStack} / \texttt{redoStack} (capped at 200 entries) stored
on \texttt{GameSessionState}. Each entry is an \texttt{EditorSnapshot}
--- a \texttt{structuredClone} of \texttt{pgnModel}, \texttt{pgnText},
\texttt{currentPly}, \texttt{selectedMoveId}, and
\texttt{pgnLayoutMode}. Undo and redo restore the snapshot and call
\texttt{onSyncChessParseState} + \texttt{onPgnChange} +
\texttt{onNavigationChange} + \texttt{onUndoRedoDepthChange}; they do
not go through \texttt{applyPgnModelUpdate} so they do not push further
history entries.

\subsubsection{7.5 The Metadata System}\label{the-metadata-system}

Metadata in X2Chess is a multi-layer system that begins with the PGN
text and ends at search-capable, typed, user-extensible field
definitions.

\paragraph{PGN headers as source of
truth}\label{pgn-headers-as-source-of-truth}

Every game carries its metadata inside its own PGN text as bracket
header lines:

\begin{verbatim}
[White "Kasparov, G."]
[Result "1-0"]
[WhiteElo "2812"]
[Date "1997.05.11"]
\end{verbatim}

The PGN text is always the authoritative record. Derived representations
(the DB index, the viewer display) are computed from it and can be
rebuilt at any time.

\paragraph{Two extraction modes}\label{two-extraction-modes}

\texttt{resource/domain/metadata.ts} exposes two extraction functions
that parse headers from raw PGN text:

\begin{itemize}
\tightlist
\item
  \textbf{\texttt{extractPgnMetadata(pgnText,\ keys?)}} --- projects the
  requested header keys as plain strings. Used where only display or
  storage is needed (listing, file/directory adapters).
\item
  \textbf{\texttt{extractHybridPgnMetadata(pgnText,\ keys?)}} ---
  projects known keys through typed parsers from
  \texttt{PGN\_METADATA\_SCHEMA}, yielding \texttt{HybridPgnMetadata}.
  Known types:
\end{itemize}

\begin{longtable}[]{@{}ll@{}}
\toprule
Key & Parsed type\tabularnewline
\midrule
\endhead
\texttt{Date} & \texttt{PgnDateValue}
(\texttt{\{\ raw,\ year,\ month,\ day\ \}}) --- \texttt{?} fields become
\texttt{null}\tabularnewline
\texttt{WhiteElo}, \texttt{BlackElo} & \texttt{number} (via
\texttt{parseInt}; non-numeric → \texttt{undefined})\tabularnewline
\texttt{Result} & \texttt{PgnResultValue} union
(\texttt{"1-0"\ \textbackslash{}\textbar{}\ "0-1"\ \textbackslash{}\textbar{}\ "1/2-1/2"\ \textbackslash{}\textbar{}\ "*"})\tabularnewline
\texttt{X2Style} & \texttt{X2StyleValue} union
(\texttt{"plain"\ \textbackslash{}\textbar{}\ "text"\ \textbackslash{}\textbar{}\ "tree"})\tabularnewline
All other known keys & \texttt{string}\tabularnewline
Unknown keys & \texttt{string} (pass-through)\tabularnewline
\bottomrule
\end{longtable}

Unknown keys are preserved as strings; the \texttt{HybridPgnMetadata}
type is the intersection of \texttt{PgnMetadataKnownValues} and
\texttt{Record\textless{}string,\ PgnMetadataScalar\ \textbar{}\ undefined\textgreater{}}.

\paragraph{Key lists}\label{key-lists}

Three ordered key lists define projection scopes:

\begin{longtable}[]{@{}ll@{}}
\toprule
\begin{minipage}[b]{0.05\columnwidth}\raggedright\strut
Constant\strut
\end{minipage} & \begin{minipage}[b]{0.05\columnwidth}\raggedright\strut
Content\strut
\end{minipage}\tabularnewline
\midrule
\endhead
\begin{minipage}[t]{0.05\columnwidth}\raggedright\strut
\texttt{PGN\_STANDARD\_METADATA\_KEYS}\strut
\end{minipage} & \begin{minipage}[t]{0.05\columnwidth}\raggedright\strut
14 keys --- the Seven Tag Roster plus ECO, Opening, WhiteElo, BlackElo,
TimeControl, Termination, Annotator\strut
\end{minipage}\tabularnewline
\begin{minipage}[t]{0.05\columnwidth}\raggedright\strut
\texttt{KNOWN\_PGN\_METADATA\_KEYS}\strut
\end{minipage} & \begin{minipage}[t]{0.05\columnwidth}\raggedright\strut
Standard keys plus \texttt{X2Style} and \texttt{Material}\strut
\end{minipage}\tabularnewline
\begin{minipage}[t]{0.05\columnwidth}\raggedright\strut
\texttt{DEFAULT\_RESOURCE\_VIEWER\_METADATA\_KEYS}\strut
\end{minipage} & \begin{minipage}[t]{0.05\columnwidth}\raggedright\strut
7-key default column set for the resource viewer: White, Black, Date,
Event, ECO, Opening, Result\strut
\end{minipage}\tabularnewline
\bottomrule
\end{longtable}

The \texttt{METADATA\_KEY} const object mirrors
\texttt{PgnMetadataKnownValues} as string constants so that renaming a
key in the type produces a compile error at every call site.

\paragraph{DB-level metadata index}\label{db-level-metadata-index}

When a game is created or saved in a \texttt{db} resource,
\texttt{writeMetadata} in \texttt{resource/adapters/db/db\_indexer.ts}
denormalizes header values into the \texttt{game\_metadata} table (one
row per key--value pair, keyed by \texttt{(game\_id,\ meta\_key)}). This
index drives:

\begin{itemize}
\tightlist
\item
  \textbf{List display} --- \texttt{db\_adapter.ts} batch-loads metadata
  for all games in a single query at list time, avoiding per-game reads.
\item
  \textbf{Text search} --- \texttt{searchByText} queries
  \texttt{game\_metadata} for substring matches on \texttt{White},
  \texttt{Black}, \texttt{Event}, and \texttt{Site}.
\end{itemize}

The \texttt{metadata\_keys} table records which keys are present in a
given database; the adapter returns this list alongside each
\texttt{list} result so the viewer knows which columns to offer.

Every save deletes and rebuilds the metadata rows for the affected game
--- the index is always derived from the current PGN text.

\paragraph{\texorpdfstring{Derived \texttt{Material}
key}{Derived Material key}}\label{derived-material-key}

For position games (\texttt{{[}SetUp\ "1"{]}} present),
\texttt{writeMetadata} also derives a \texttt{Material} key by passing
the \texttt{{[}FEN{]}} header value through \texttt{materialKeyFromFen}
in \texttt{resource/domain/material\_key.ts}. The key has the form
\texttt{K\textless{}white-pieces\textgreater{}vK\textless{}black-pieces\textgreater{}}
with pieces sorted by descending value, e.g. \texttt{KQPPPvKRP},
\texttt{KBNvK}, \texttt{KvK}. It enables substring search for endgame
material configurations.

\paragraph{User-defined schemas}\label{user-defined-schemas}

Beyond the standard PGN keys, users can define custom metadata schemas.
A \texttt{MetadataSchema} contains:

\begin{itemize}
\tightlist
\item
  A stable UUID \texttt{id} and a display \texttt{name}.
\item
  A monotonically incrementing \texttt{version} (bumped on each save).
\item
  An ordered list of \texttt{MetadataFieldDefinition} entries, each with
  a \texttt{key}, \texttt{label}, \texttt{type}, \texttt{required} flag,
  \texttt{orderIndex} (10, 20, 30, \ldots{} with gaps for future
  insertions), and optional \texttt{selectValues} and
  \texttt{description}.
\end{itemize}

\textbf{Field types:} \texttt{text}, \texttt{date}, \texttt{select},
\texttt{number}, \texttt{flag}, \texttt{game\_link}.

Schemas are persisted in localStorage under
\texttt{x2chess.metadata-schemas} via the versioned store in
\texttt{frontend/src/resources\_viewer/schema\_storage.ts}. They can be
exported and imported as JSON files (format marker:
\texttt{"x2chess-schema":\ "1"}).

The built-in \texttt{BUILT\_IN\_SCHEMA} (id \texttt{"builtin"}) is a
read-only schema for the standard PGN Seven Tag Roster plus common
extensions. It is used when no user-defined schema is associated with a
resource.

\paragraph{Relationship between
layers}\label{relationship-between-layers}

\begin{verbatim}
PGN text (source of truth)
  │
  ├─ extractPgnMetadata()       → plain string map
  │    used by: file/directory adapters, db list display, db indexer
  │
  ├─ extractHybridPgnMetadata() → typed HybridPgnMetadata
  │    used by: resource viewer columns, typed field rendering
  │
  └─ writeMetadata() [db only]  → game_metadata + metadata_keys rows
       used by: text search, batch list loading

MetadataSchema (localStorage)
  └─ maps field keys to types and labels
       used by: resource viewer column configuration, metadata editor dialog
\end{verbatim}

\subsection{8) Constraints for Ongoing
Work}\label{constraints-for-ongoing-work}

\begin{itemize}
\tightlist
\item
  Keep browser fat-client as primary source of feature truth.
\item
  Keep reduced app-profile differences explicit and documented.
\item
  Keep Tauri as reduced app shell, not a replacement for frontend domain
  logic.
\item
  Keep local resource APIs explicit and maintainable.
\item
  Keep local resource read/write in client code only (no backend
  mediation).
\item
  Keep remote integrations behind adapter boundaries and contracts.
\item
  Keep top-level \texttt{resource/} as canonical home for PGN resource
  contracts, kind routing, and metadata schema.
\item
  Keep canonical resource kinds explicit as \texttt{file},
  \texttt{directory}, and \texttt{db}.
\item
  Keep SQL schema/migration ownership in \texttt{resource/database/}
  (never under frontend domain modules).
\item
  Keep backend optional and decoupled from local resource workflows.
\item
  Enforce Component-Contract on components (integration + configuration
  + communication APIs, documented publicly in code).
\item
  Enforce clean/professional CSS defaults and maintain explicit theming
  extension points for future user adaptation.
\item
  Keep local runtime data scenario-separated under \texttt{run/}.
\item
  Keep production runtime data outside repository paths (explicit prod
  root or OS app-data default).
\item
  Keep scenario game libraries under
  \texttt{run/\textless{}SCENARIO\textgreater{}/games} for local
  scenarios.
\item
  Keep runtime configuration layered: default app config + partial user
  override in data area.
\item
  React \texttt{useReducer} (\texttt{AppStoreState}) is the
  authoritative state for rendering; all component reads go through
  selectors, all component writes go through \texttt{ServiceContext}.
\item
  \texttt{activeSessionRef.current} (\texttt{GameSessionState}) is a
  deliberate mutable working store for the active session --- services
  write to it during operations and flush to React via callbacks or
  \texttt{flushSessionState()} when the operation completes (see §7.3).
  This is an intentional performance trade-off, not a violation.
\item
  Services own their internal state as closure variables; no shared
  mutable \texttt{AppState} object.
\item
  Keep service initialisation in \texttt{useAppStartup}; expose service
  callbacks exclusively via \texttt{ServiceContext}.
\item
  Pure-logic modules (\texttt{model/}, \texttt{editor/},
  \texttt{board/}, \texttt{game\_sessions/}, \texttt{resources/}) must
  remain framework-free.
\item
  Update this manual when architecture decisions, boundaries, or
  delivery paths change.
\end{itemize}

\end{document}
